% GENERATED FILE. Edit canonical lessons, then run npm run generate:books.
% canonical-source-hash: fnv1a64:c612fd735baf1ddc
% canonical-lessons: UR-C415-dhaba, UR-C415-thela, UR-C415-pan, UR-C415-zaiqa, UR-C415-parosna

\chapter{Street Food}
\label{ch:ur-street-food}
\hlchaptermodality{415}

\begin{chapteropening}
\textbf{By the end of this chapter:} \emph{I can eat at a roadside stall and say how the food tastes.}

Complete the last lesson of chapter 415: I can eat at a roadside stall and say how the food tastes.
\end{chapteropening}

\section[\texorpdfstring{\ur{ڈھابا} (ḍhābā)}{ڈھابا (ḍhābā)}]{\ur{ڈھابا} (ḍhābā) — a roadside eatery}

\label{lesson:UR-C415-dhaba}



\begin{quote}
Before the new one: say the Urdu for to growl, then the Urdu for to chase off.
\end{quote}

\subsection*{You'll want to know: \ur{ڈھابا}}
\textbf{\ur{ڈھابا}} — \emph{ḍhābā} — \textquotedblleft{}a roadside eatery\textquotedblright{}.

\begin{quote}
\ur{شاہراہ} \ur{کے} \ur{کنارے} \ur{ایک} \ur{ڈھابا} \ur{ہے۔}
\end{quote}

\begin{quote}
\emph{shāhrāh ke kināre ek ḍhābā hai.}
\end{quote}

\begin{quote}
There's a roadside eatery beside the highway.
\end{quote}

Masculine.

\begin{cousinweb}
\textbf{\ur{ڈھابا}} (\emph{ḍhābā}) is inherited: CDIAL lists Hindi \emph{ḍhābā}, \textquotedblleft{}thatched roofing\textquotedblright{}, under \emph{ḍhapp}, \textquotedblleft{}cover\textquotedblright{}, a form it reconstructs (CDIAL 5579); the sense \textquotedblleft{}a roadside eatery\textquotedblright{}, a place under a thatch, is probably a later one.
\end{cousinweb}

\begin{grammarlens}[title={What you've built}]
One more word for getting around.
\end{grammarlens}

\subsection*{Your turn}
\begin{itemize}
\raggedright
  \item \emph{Say it:} \emph{ḍhābā}
  \item \emph{Say it:} \emph{ḍhābā}, once more
  \item \emph{Say it:} \emph{shāhrāh ke kināre ek ḍhābā hai.} — There's a roadside eatery beside the highway.
  \item \emph{Recall:} say the Urdu for to growl, then the Urdu for to chase off, then say \emph{ḍhābā} again
\end{itemize}

\begin{tcolorbox}[breakable,colback=teal!4,colframe=teal!35!black,title={Before you move on}]
What does \ur{ڈھابا} mean? (\textquotedblleft{}A roadside eatery\textquotedblright{}.) Say it once more.
\end{tcolorbox}
\section[\texorpdfstring{\ur{ٹھیلا} (ṭhelā)}{ٹھیلا (ṭhelā)}]{\ur{ٹھیلا} (ṭhelā) — a handcart}

\label{lesson:UR-C415-thela}



\begin{quote}
Before the new one: say the Urdu for to chase off, then the Urdu for a roadside eatery.
\end{quote}

\subsection*{You'll want to know: \ur{ٹھیلا}}
\textbf{\ur{ٹھیلا}} — \emph{ṭhelā} — \textquotedblleft{}a handcart\textquotedblright{}.

\begin{quote}
\ur{ٹھیلے} \ur{پر} \ur{ایک} \ur{آدمی} \ur{پھل} \ur{بیچتا} \ur{ہے۔}
\end{quote}

\begin{quote}
\emph{ṭhele par ek ādmī phal bechtā hai.}
\end{quote}

\begin{quote}
A man sells fruit from a handcart.
\end{quote}

Masculine.

\begin{cousinweb}
\textbf{\ur{ٹھیلا}} (\emph{ṭhelā}) is inherited: CDIAL lists the Hindi cousin \emph{ṭhelā}, \textquotedblleft{}a push-cart\textquotedblright{}, under \emph{ṭhēḍḍ}, \textquotedblleft{}push\textquotedblright{}, a form it reconstructs (CDIAL 5512).
\end{cousinweb}

\begin{grammarlens}[title={What you've built}]
One more word for getting around.
\end{grammarlens}

\subsection*{Your turn}
\begin{itemize}
\raggedright
  \item \emph{Say it:} \emph{ṭhelā}
  \item \emph{Say it:} \emph{ṭhelā}, once more
  \item \emph{Say it:} \emph{ṭhele par ek ādmī phal bechtā hai.} — A man sells fruit from a handcart.
  \item \emph{Recall:} say the Urdu for to chase off, then the Urdu for a roadside eatery, then say \emph{ṭhelā} again
\end{itemize}

\begin{tcolorbox}[breakable,colback=teal!4,colframe=teal!35!black,title={Before you move on}]
What does \ur{ٹھیلا} mean? (\textquotedblleft{}A handcart\textquotedblright{}.) Say it once more.
\end{tcolorbox}
\section[\texorpdfstring{\ur{پان} (pān)}{پان (pān)}]{\ur{پان} (pān) — a betel leaf; paan}

\label{lesson:UR-C415-pan}



\begin{quote}
Before the new one: say the Urdu for a roadside eatery, then the Urdu for a handcart.
\end{quote}

\subsection*{You'll want to know: \ur{پان}}
\textbf{\ur{پان}} — \emph{pān} — \textquotedblleft{}a betel leaf; paan\textquotedblright{}.

\begin{quote}
\ur{کھانے} \ur{کے} \ur{بعد} \ur{کچھ} \ur{لوگ} \ur{پان} \ur{کھاتے} \ur{ہیں۔}
\end{quote}

\begin{quote}
\emph{khāne ke bād kuchh log pān khāte haiṅ.}
\end{quote}

\begin{quote}
After a meal some people chew paan.
\end{quote}

Masculine.

\begin{cousinweb}
\textbf{\ur{پان}} (\emph{pān}) is inherited: CDIAL lists the Hindi cousin \emph{pān}, \textquotedblleft{}a betel leaf\textquotedblright{}, under Sanskrit \emph{parṇá}, \textquotedblleft{}a leaf\textquotedblright{} (CDIAL 7918).
\end{cousinweb}

\begin{grammarlens}[title={What you've built}]
One more word for getting around.
\end{grammarlens}

\subsection*{Your turn}
\begin{itemize}
\raggedright
  \item \emph{Say it:} \emph{pān}
  \item \emph{Say it:} \emph{pān}, once more
  \item \emph{Say it:} \emph{khāne ke bād kuchh log pān khāte haiṅ.} — After a meal some people chew paan.
  \item \emph{Recall:} say the Urdu for a roadside eatery, then the Urdu for a handcart, then say \emph{pān} again
\end{itemize}

\begin{tcolorbox}[breakable,colback=teal!4,colframe=teal!35!black,title={Before you move on}]
What does \ur{پان} mean? (\textquotedblleft{}A betel leaf; paan\textquotedblright{}.) Say it once more.
\end{tcolorbox}
\section[\texorpdfstring{\ur{ذائقہ} (zāiqā)}{ذائقہ (zāiqā)}]{\ur{ذائقہ} (zāiqā) — a flavour; a taste}

\label{lesson:UR-C415-zaiqa}



\begin{quote}
Before the new one: say the Urdu for a handcart, then the Urdu for a betel leaf.
\end{quote}

\subsection*{You'll want to know: \ur{ذائقہ}}
\textbf{\ur{ذائقہ}} — \emph{zāiqā} — \textquotedblleft{}a flavour; a taste\textquotedblright{}.

\begin{quote}
\ur{اس} \ur{ڈھابے} \ur{کے} \ur{کھانے} \ur{کا} \ur{ذائقہ} \ur{اچھا} \ur{ہے۔}
\end{quote}

\begin{quote}
\emph{is ḍhābe ke khāne kā zāiqā achchhā hai.}
\end{quote}

\begin{quote}
The food at this eatery tastes good.
\end{quote}

Masculine.

\begin{cousinweb}
is usually traced, through Persian, to Arabic \emph{dhāʾiqa}, \textquotedblleft{}the sense of taste\textquotedblright{}; a dish with a good \textbf{\ur{ذائقہ}} is \textbf{\ur{مزیدار}}.
\end{cousinweb}

\begin{grammarlens}[title={What you've built}]
One more word for getting around.
\end{grammarlens}

\subsection*{Your turn}
\begin{itemize}
\raggedright
  \item \emph{Say it:} \emph{zāiqā}
  \item \emph{Say it:} \emph{zāiqā}, once more
  \item \emph{Say it:} \emph{is ḍhābe ke khāne kā zāiqā achchhā hai.} — The food at this eatery tastes good.
  \item \emph{Recall:} say the Urdu for a handcart, then the Urdu for a betel leaf, then say \emph{zāiqā} again
\end{itemize}

\begin{tcolorbox}[breakable,colback=teal!4,colframe=teal!35!black,title={Before you move on}]
What does \ur{ذائقہ} mean? (\textquotedblleft{}A flavour; a taste\textquotedblright{}.) Say it once more.
\end{tcolorbox}
\section[\texorpdfstring{\ur{پروسنا} (parosnā)}{پروسنا (parosnā)}]{\ur{پروسنا} (parosnā) — to serve (food)}

\label{lesson:UR-C415-parosna}



\begin{quote}
Before the new one: say the Urdu for a betel leaf, then the Urdu for a flavour.
\end{quote}

\subsection*{You'll want to know: \ur{پروسنا}}
\textbf{\ur{پروسنا}} — \emph{parosnā} — \textquotedblleft{}to serve (food)\textquotedblright{}.

\begin{quote}
\ur{ڈھابے} \ur{والے} \ur{نے} \ur{ہمیں} \ur{گرم} \ur{روٹی} \ur{پروسی۔}
\end{quote}

\begin{quote}
\emph{ḍhābe vāle ne hameṅ garm roṭī parosī.}
\end{quote}

\begin{quote}
The man at the eatery served us hot bread.
\end{quote}

Past, with \ur{نے} (he or she …): \ur{اس} \ur{نے} \ur{کھانا} \ur{پروسا}, \emph{us ne khānā parosā}.

\begin{cousinweb}
\textbf{\ur{پروسنا}} (\emph{parosnā}) is inherited: CDIAL lists the Hindi cousin \emph{parosnā}, \textquotedblleft{}to offer food to a guest\textquotedblright{}, under Sanskrit \emph{párivēṣati} (CDIAL 7888).
\end{cousinweb}

\begin{grammarlens}[title={What you've built}]
One more word for getting around.
\end{grammarlens}

\subsection*{Your turn}
\begin{itemize}
\raggedright
  \item \emph{Say it:} \emph{parosnā}
  \item \emph{Say it:} \emph{parosnā}, once more
  \item \emph{Say it:} \emph{ḍhābe vāle ne hameṅ garm roṭī parosī.} — The man at the eatery served us hot bread.
  \item \emph{Recall:} say the Urdu for a betel leaf, then the Urdu for a flavour, then say \emph{parosnā} again
\end{itemize}

\begin{tcolorbox}[breakable,colback=teal!4,colframe=teal!35!black,title={Before you move on}]
What does \ur{پروسنا} mean? (\textquotedblleft{}To serve (food)\textquotedblright{}.) Say it once more.
\end{tcolorbox}
